\documentclass[11pt]{article}

\usepackage[utf8]{inputenc}
\usepackage[T1]{fontenc}
\usepackage{amsmath}
\usepackage{amssymb}
\usepackage{booktabs}
\usepackage{graphicx}
\usepackage[margin=1.1in]{geometry}
\usepackage[colorlinks=true,linkcolor=blue,citecolor=blue]{hyperref}

\title{Cache-Aware Scheduling for Streaming Joins}
\author{A. Author \and B. Author}
\date{June 2026}

\begin{document}

\maketitle

\begin{abstract}
Streaming join operators spend most of their time waiting on memory, not
on comparisons. We describe a scheduler that orders probe batches by the
cache residency of the hash partitions they touch, and show that the
reordering recovers most of the throughput a partition-oblivious
scheduler leaves behind.
\end{abstract}

\section{Introduction}
\label{sec:intro}

A streaming hash join keeps one side of the join resident as a partitioned
hash table and probes it with batches arriving from the other side. The
textbook cost model counts comparisons, which suggests that the operator's
throughput should be flat in the number of partitions. Measured on real
hardware it is not: throughput falls by more than a factor of two once the
resident side outgrows the last-level cache, and the fall is entirely
attributable to partitions being evicted between the probes that need
them~\cite{balkesen2013}.

The scheduler is free to choose the order in which buffered probe batches
are dispatched. A partition-oblivious scheduler dispatches in arrival
order, which interleaves partitions arbitrarily and evicts each one
shortly before it is needed again. We show that a scheduler which groups
buffered batches by target partition, and dispatches whole groups,
recovers most of that loss.

Our contributions are:

\begin{itemize}
\item A residency model that predicts, from partition sizes and the
  buffered batch mix alone, how much reordering can recover
  (Section~\ref{sec:method}).
\item A scheduler that applies the model online, with a bounded buffer
  and no change to output order.
\item An evaluation on three join workloads showing a
  \(1.7\times\) mean throughput gain (Section~\ref{sec:results}).
\end{itemize}

The rest of this paper is organized around those three pieces.


\section{Method}
\label{sec:method}

\subsection{A residency model}

Let \(P\) be the set of hash partitions and \(s_p\) the resident size of
partition \(p \in P\). A probe batch \(b\) targets exactly one partition,
written \(\pi(b)\). For a dispatch order \(\sigma\) over a buffer of \(n\)
batches, the reuse distance of the \(i\)-th batch is the total resident
size touched since the previous batch with the same target:

\begin{equation}
\label{eq:reuse}
d_\sigma(i) = \sum_{j \in R_\sigma(i)} s_{\pi(\sigma_j)},
\qquad
R_\sigma(i) = \{\, j < i : \pi(\sigma_j) \neq \pi(\sigma_i) \,\}.
\end{equation}

A batch hits cache when \(d_\sigma(i) < C\), the effective cache capacity.
Minimizing misses is therefore minimizing the number of indices whose
reuse distance exceeds \(C\), which grouping by target achieves exactly:
consecutive batches with equal targets have reuse distance zero.

\subsection{The scheduler}

The online scheduler holds a bounded buffer of \(n\) batches and always
dispatches the group with the largest count, breaking ties by arrival
time so that no batch can be starved. The bound on \(n\) is what keeps
latency predictable:

\begin{itemize}
\item A batch waits at most \(n - 1\) dispatch slots.
\item Output order is restored by the operator's existing sequence
  numbers, so no downstream operator observes the reordering.
\end{itemize}

\begin{table}[htbp]
\centering
\begin{tabular}{lrr}
\toprule
Buffer \(n\) & Miss rate & Throughput \\
\midrule
1 (arrival order) & 0.41 & 1.00 \\
8 & 0.22 & 1.44 \\
32 & 0.13 & 1.71 \\
128 & 0.11 & 1.74 \\
\bottomrule
\end{tabular}
\caption{Buffer size against measured miss rate and normalized
throughput on the \textsc{tpch-q3} workload.}
\label{tab:buffer}
\end{table}

Table~\ref{tab:buffer} shows that the returns flatten by \(n = 32\),
which is the value used throughout Section~\ref{sec:results}.


\section{Results}
\label{sec:results}

We evaluate on three workloads: \textsc{tpch-q3}, a synthetic uniform
join, and a skewed join with a Zipfian key distribution. All runs use a
single core with a 32~MiB last-level cache, and all numbers are the median
of eleven runs.

\begin{table}[htbp]
\centering
\begin{tabular}{lrrr}
\toprule
Workload & Arrival order & Grouped & Gain \\
\midrule
\textsc{tpch-q3} & 1.00 & 1.71 & \(1.71\times\) \\
Uniform & 1.00 & 1.66 & \(1.66\times\) \\
Zipfian & 1.00 & 1.79 & \(1.79\times\) \\
\bottomrule
\end{tabular}
\caption{Normalized throughput by workload, buffer \(n = 32\).}
\label{tab:main}
\end{table}

The gain is largest on the skewed workload, which is the case the
residency model predicts: a skewed key distribution concentrates probes
on few partitions, so grouping keeps those partitions resident across
many more batches than a uniform mix does.

\subsection{Where the model is wrong}

The model assumes a partition is either resident or absent. Real caches
hold partial partitions, so the predicted miss rate is pessimistic by
roughly four points across every workload we measured; the predicted
\emph{ordering} of configurations was correct in every case, which is
what the scheduler actually consumes~\cite{blanas2011}.


\section{Conclusion}
\label{sec:conclusion}

Reordering probe batches by partition residency is a scheduling change,
not a data-structure change: it needs no extra index, no additional
memory, and no change to the join's output order. That is what makes it
worth doing.

\bibliographystyle{plain}
\bibliography{refs}

\end{document}
